% Further NN-G036 refinement; root whole-proof PASS 2026-11-22T00:12:35Z.
\paragraph{A mass-six action gives integral divisibility.}
Suppose $W\in\mathbb Z[D]$ satisfies $W\equiv D\pmod2$
and $w=W(4)$ is divisible by four.
Write $W=D+2R$. Then $R(4)=(w-4)/2$ is even,
and $R(-2)\equiv R(4)\pmod2$. Therefore
\[
                         W(-2)\equiv2\pmod4.        \tag{G036.U1}
\]
In particular $W$ and $D+2$ are coprime over
$\mathbb Q[D]$.
If $(D+2)G=W F$ for integral polynomials, then
$W\mid G$ in $\mathbb Z[D]$, even when $W$ is
not primitive. To check the content explicitly,
write $W=cW_0$ with $W_0$ primitive and $c>0$.
Gauss's content identity and primitivity of $D+2$
first give $c\mid G$. Write $G=cG_0$.
Coprimality gives $W_0\mid G_0$ over $\mathbb Q[D]$,
and primitivity makes this quotient integral.

Consequently, in any of the closed tables under
consideration, $u_i=6$ forces $U_i=D+2$, and
\[
                 4\mid w,\ u_i=6
                 \quad\Longrightarrow\quad w\mid g_j
                 \qquad(i\ne j).                  \tag{G036.U2}
\]
At $w=12$ this excludes both $(g_1,g_2)=(72,112)$
rows, since $u_1=6$ but $12\nmid112$.
It also excludes the three $(80,96)$ rows with
$d_2=14,u_2=6$, and their three index-reversed rows:
in each case the opposite action has mass eighty.
Thus eight entire base rows, with all their auxiliary
increments, are impossible.

\paragraph{The forward row with action masses eighty and ninety-six.}
Consider
\[
 (w,g_1,g_2,d_1,d_2,u_1,u_2)
                    =(12,80,96,12,4,4,12).         \tag{G036.U3}
\]
Here $U_1=D$ and $\delta_2=\psi$.
If $W(0)\ne0$, the integral quotient argument for
$U_1=D$ gives
$G_2=DWC$ with $C\equiv D-1\pmod2$.
It applies to nonprimitive $W$ as well by the same
content argument just proved. Its value would be
$C(4)=96/(4\cdot12)=2$, whereas its displayed
parity requires an odd value. This is impossible.

It remains that $W=DH$, with the mass-twelve
classification proved above: $H(4)=3$,
$H\equiv1\pmod2$, and $H(0)\ne\pm1$.
Either $H=3$, or $H$ is primitive irreducible
of leading coefficient two.
Since three does not divide $g_1=80$, in the
nonconstant case $H\nmid G_1$.
The second cell identity then forces $H\mid U_2$.
As in the earlier mass-twelve proof, $H(0)$ is
odd and not $\pm1$, whereas $U_2(0)$ is zero
or $\pm8$; therefore $U_2(0)=0$.
Writing $U_2=DJ$, both $H,J$ are primitive
irreducible of leading coefficient two and mass
three, so $J=H$ and $U_2=W$.
For $H=3$, reducing the same identity modulo
three instead gives $3\mid U_2$ coefficientwise;
its effective mass twelve and parity force
$U_2=3D=W$. Thus $U_2=W$ in every case.

The two cell identities now imply
\[
 G_2=G_1+D\psi,\qquad
                    (H-1)G_1=D(H\delta_1+\psi).     \tag{G036.U4}
\]
Write $\delta_1=\psi+2R$, where $R$ is effective
of mass four. Then $\delta_1(0)\in\{-8,0,8\}$.
If $\delta_1(0)=0$, both $\delta_1$ and $\psi$
are divisible by $D$. Since $H(0)\ne1$,
(G036.U4) forces $D^2\mid G_1$.
The quotient has even coefficients but value
$80/16=5$ at four, a contradiction.

If $\delta_1(0)=\pm8$, first $D\mid G_1$.
Divide (G036.U4) by $D$ and evaluate at zero.
Coprimality of the consecutive integers
$H(0),H(0)-1$ shows that $H(0)-1$ divides eight.
The explicit root classification gives
\[
 H(0)=1+2o^2\quad\hbox{or}\quad
 H(0)=1+2(o^2-e^2),
\]
with $o$ positive odd and $e$ positive even.
In particular $H(0)\equiv3\pmod8$.
The only such odd integer with $H(0)-1\mid8$
is three. The first formula then forces $o=1$;
the second would require $o^2-e^2=1$,
impossible for positive $e$.
Consequently $H=3$ and $W=3D$.

Equation (G036.U4) becomes
$G_1=D(2\psi+3R)$.
Its even parity forces $R=2P$ integrally, where
$P$ is effective of mass two, hence either two
or some $T_k$. Thus
\[
 \delta_1=\psi+4P,\qquad
 G_1=D(2\psi+6P),\qquad
 G_2=D(3\psi+6P).
\]
The central Laurent coefficient of $2D\psi$
is four. For $P=T_k$ with $k\ge2$, that of
$6DP$ is zero, whereas $G_1$ must have central
coefficient at least sixteen.
Therefore $P=2$ or $T_1$.

Put $L=2\psi+6P$, an integral quadratic with
leading coefficient two and $L(4)=20$.
The second auxiliary identity gives
$V_2L=3(DG_2-A_2)$.
Modulo three, $L$ reduces to $2\psi\ne0$,
so $V_2=3Z$ integrally.
Since $V_2$ is effective and $V_2\equiv D\pmod2$,
$Z$ is effective and $Z\equiv D\pmod2$.
The actual second auxiliary coefficient is
\[
              A_2=D^2(3\psi+6P)-Z(2\psi+6P).       \tag{G036.U5}
\]
The first term has degree four and leading
coefficient three. A $Z$ of degree greater
than two makes the leading coefficient of
$A_2$ negative. If $\deg Z=2$, its leading
coefficient is positive even and at least two,
so the degree-four coefficient is at most
$3-2\cdot2<0$.
Thus effectivity forces $\deg Z\le1$.
But then $V_2=3Z$ has no $q^2$ coefficient,
contradicting its required coefficient at least two.
The forward row (G036.U3) is excluded.
\hfill$\square$

\paragraph{Geometric scope.}
The divisibility lemma removes ordinary-dimension
patterns for effective graded sign-representation
objects in the specified closed kernel presentation.
It does not assume positive polynomial quotients:
content and integrality alone supply the character
divisibilities. The forward-row exclusion further
uses actual coefficientwise effectivity of the
structure objects. These are the same mathematical
objects of NN-G010 and SGA~3, I, 4.7.2--4.7.3,
with composition in the supplied archive
cote~2, scan~17 transcription. No change of
multiplication-table presentation is permitted
in this conclusion.
